% SPDX-License-Identifier: CC-BY-SA-4.0
% Author: Matthieu Perrin
% Part: <Nom de la partie>
% Section: <Nom de la section>
% Sub-section: <Nom de la sous-section>  % (facultatif, laisser vide si non utilisé)
% Frame: <Titre de la slide>

\begingroup

\begin{frame}{Théorème d'incomplétude de Gödel}

    \vspace{-2mm}
  \begin{block}{Théorème -- Gödel 1931 (admis)}
    Il n'existe pas de système de preuve $\langle \Phi, \Pi, \valid \rangle$ à la fois : 
    \begin{description}[Peano expressif :]
    \item[Formel :] la vérification des preuves est mécanisable
      \begin{itemize}
      \item $\textsc{Entscheidungsproblem} \in \textsc{re}$
        \footnote{Ou, de manière équivalente, $\textsc{Vérification des preuves} \in \textsc{r}$}
      \end{itemize}
    \item[Peano expressif :] possibilité d'exprimer des questions sur les nombre entiers
      \begin{itemize}
      \item Les axiomes de Dedekind-Peano sont démontrables\footnote{$\exists \mathbb{N},\exists 0 \in \mathbb{N}, \exists s :\mathbb{N} \rightarrow {\mathbb{N}\setminus \{0\}} \text{ injective}, \forall P\subseteq \mathbb{N}, (0 \in P \land s(P) \subseteq P) \rightarrow P=\mathbb{N}$}
      \end{itemize}
    \item[Cohérent :] on ne peut pas montrer une proposition et son opposée
      \begin{itemize}
      \item $\forall \varphi\in \Phi,\quad(\nexists \pi, \pi \valid \varphi) \quad \lor\quad (\nexists \pi, \pi \valid \lnot\varphi)$
      \end{itemize}
    \item[Complet :] toute proposition peut être démontrée ou invalidée
      \begin{itemize}
      \item $\forall \varphi\in \Phi,\quad (\exists \pi, \pi \valid \varphi) \quad \lor\quad (\exists \pi, \pi \valid \lnot\varphi)$
      \end{itemize}
    \end{description}
  \end{block}
  
  \pause
    \vspace{-2mm}
  \begin{exampleblock}{Exemple -- Hypothèse du Continu}
    Cette proposition est indécidable dans les mathématiques modernes (ZFC) :
    \vspace{-1mm}
    $$\example{\exists E,\quad \mathbb{N} \subset E \subset \mathbb{R} \quad\land\quad |\mathbb{N}| < |E| < |\mathbb{R}|}$$ 

    \vspace{-3mm}
  \end{exampleblock}
  
\end{frame}

\endgroup
